\section{Lecture 21: Convolutional Neural Networks}

\subsection{Motivation}

So far, we have studied neural networks as general function approximators.

\medskip

\noindent
However, many real-world data types exhibit structure. In particular:

\begin{itemize}
    \item images have spatial structure,
    \item nearby pixels are highly correlated,
    \item patterns repeat across locations.
\end{itemize}

\medskip

\noindent
\textbf{Key question.}  
How can we design neural networks that exploit such structure?

\medskip

\noindent
\textbf{Guiding idea.}  
Convolutional neural networks (CNNs) incorporate prior knowledge about spatial locality and translation invariance.

\subsection{Convolution as a Structured Linear Operator}

A convolutional layer applies a filter (kernel) across the input.

\medskip

\noindent
For a 1D signal:
\[
y[i] = \sum_{k} w[k] \, x[i-k].
\]

\medskip

\noindent
For images (2D):
\[
y[i,j] = \sum_{u,v} w[u,v] \, x[i-u, j-v].
\]

\medskip

\noindent
\textbf{Interpretation.}

\begin{itemize}
    \item Local connectivity: each output depends on a small neighborhood
    \item Weight sharing: the same filter is applied everywhere
\end{itemize}

\medskip

\noindent
\textbf{Matrix view.}

\begin{itemize}
    \item Convolution corresponds to a structured (Toeplitz) matrix
    \item Far fewer parameters than fully connected layers
\end{itemize}

\medskip

\noindent
\textbf{Insight.}  
Convolution imposes strong structure on linear transformations.

\subsection{Translation Invariance}

A key property of convolution:

\medskip

\noindent
\textbf{Translation equivariance.}

\begin{itemize}
    \item Shifting the input shifts the output
\end{itemize}

\medskip

\noindent
\textbf{Why this matters.}

\begin{itemize}
    \item Objects may appear anywhere in an image
    \item Model should detect features regardless of location
\end{itemize}

\medskip

\noindent
\textbf{Approximate invariance.}

\begin{itemize}
    \item Pooling and deeper layers help achieve invariance
\end{itemize}

\medskip

\noindent
\textbf{Insight.}  
CNNs encode translation symmetry as an inductive bias.

\subsection{Pooling and Hierarchies}

\textbf{Pooling layers.}

\begin{itemize}
    \item Reduce spatial resolution
    \item Aggregate local information
\end{itemize}

\medskip

\noindent
\textbf{Examples.}

\begin{itemize}
    \item max pooling
    \item average pooling
\end{itemize}

\medskip

\noindent
\textbf{Effect.}

\begin{itemize}
    \item increases receptive field
    \item reduces sensitivity to small shifts
\end{itemize}

\medskip

\noindent
\textbf{Hierarchical representations.}

\begin{itemize}
    \item early layers: edges, textures
    \item deeper layers: shapes, objects
\end{itemize}

\medskip

\noindent
\textbf{Insight.}  
CNNs build representations from local to global structure.

\subsection{Architectural Components}

Modern CNNs combine several elements:

\begin{itemize}
    \item convolution layers
    \item nonlinear activations (e.g., ReLU)
    \item normalization (e.g., batch normalization)
    \item pooling or strided convolution
    \item residual connections
\end{itemize}

\medskip

\noindent
\textbf{Design pattern.}

\[
\text{Conv} \rightarrow \text{BN} \rightarrow \text{ReLU}
\]

\medskip

\noindent
\textbf{Insight.}  
CNNs integrate multiple principles developed in earlier lectures.

\subsection{Applications in Vision}

CNNs are widely used in:

\begin{itemize}
    \item image classification
    \item object detection
    \item segmentation
\end{itemize}

\medskip

\noindent
\textbf{Why CNNs work well.}

\begin{itemize}
    \item exploit spatial locality
    \item reduce number of parameters
    \item encode useful invariances
\end{itemize}

\medskip

\noindent
\textbf{Limitations.}

\begin{itemize}
    \item limited ability to model long-range dependencies
    \item fixed receptive field growth
\end{itemize}

\medskip

\noindent
\textbf{Insight.}  
CNNs are powerful but tailored to specific data structures.

\subsection{Fourier Perspective on Convolution (Advanced)}

Convolution has a fundamental interpretation in the frequency domain.

\medskip

\noindent
\textbf{Fourier transform.}  
Let $\mathcal{F}$ denote the Fourier transform. For a signal $x$:
\[
\hat{x}(\omega) = \mathcal{F}[x](\omega).
\]

\medskip

\noindent
\textbf{Convolution theorem.}  
For signals $x$ and $w$:
\[
x * w \;\longleftrightarrow\; \hat{x}(\omega)\,\hat{w}(\omega),
\]
i.e., convolution in the spatial domain corresponds to pointwise multiplication in the frequency domain.

\medskip

\noindent
\textbf{Interpretation.}

\begin{itemize}
    \item Convolution acts as a \emph{filter} in frequency space
    \item The kernel $w$ determines which frequencies are amplified or suppressed
\end{itemize}

\medskip

\noindent
\textbf{Examples.}

\begin{itemize}
    \item Smooth filters $\Rightarrow$ low-pass (remove high-frequency noise)
    \item Edge detectors $\Rightarrow$ emphasize high-frequency components
\end{itemize}

\medskip

\noindent
\textbf{Insight.}  
Convolutional layers learn frequency-selective filters.

\subsection{Spectral View of CNN Layers}

Consider a convolutional layer:
\[
y = w * x.
\]

\medskip

\noindent
In the Fourier domain:
\[
\hat{y}(\omega) = \hat{w}(\omega)\,\hat{x}(\omega).
\]

\medskip

\noindent
\textbf{Implications.}

\begin{itemize}
    \item Each layer modulates the spectrum of the signal
    \item Deep CNNs apply successive spectral transformations
\end{itemize}

\medskip

\noindent
\textbf{Insight.}  
CNNs can be interpreted as hierarchical spectral filters.

\subsection{Connection to Linear Operators}

From linear algebra:

\begin{itemize}
    \item Convolution corresponds to a structured matrix (Toeplitz / circulant)
    \item Fourier transform approximately diagonalizes convolution
\end{itemize}

\medskip

\noindent
\textbf{Key idea.}

\[
\text{Convolution} \;\approx\; \text{diagonal operator in Fourier basis}.
\]

\medskip

\noindent
\textbf{Insight.}  
The Fourier basis reveals the intrinsic structure of convolutional operators.

\subsection{Role of Nonlinearity}

CNNs are not purely linear:

\begin{itemize}
    \item nonlinear activations break simple spectral interpretation
    \item interactions between frequencies emerge across layers
\end{itemize}

\medskip

\noindent
\textbf{Interpretation.}

\begin{itemize}
    \item linear part: frequency filtering
    \item nonlinear part: mixing and recombining frequencies
\end{itemize}

\medskip

\noindent
\textbf{Insight.}  
Deep CNNs combine spectral filtering with nonlinear feature interactions.

\subsection{Broader Connections}

This perspective connects CNNs to broader mathematical frameworks:

\begin{itemize}
    \item \textbf{Signal processing:} classical filtering theory
    \item \textbf{Harmonic analysis:} decomposition into frequency components
    \item \textbf{Operator theory:} diagonalization of structured operators
\end{itemize}

\medskip

\noindent
\textbf{Modern developments.}

\begin{itemize}
    \item spectral bias in neural networks
    \item frequency-domain analysis of generalization
\end{itemize}

\medskip

\noindent
\textbf{Key insight.}  
Convolutional neural networks can be understood as learnable spectral filters acting on structured data.

\subsection{A Unifying Perspective}

CNNs illustrate a central theme in machine learning:

\begin{itemize}
    \item incorporating structure into models improves efficiency and performance.
\end{itemize}

\medskip

\noindent
They combine:

\begin{itemize}
    \item inductive bias (locality, translation symmetry),
    \item deep representations,
    \item efficient parameterization.
\end{itemize}

\medskip

\noindent
\textbf{Core idea.}  
Architectural design encodes prior knowledge about the data.

\medskip

\noindent
\textbf{Key insight.}  
Good architectures reduce the burden on learning by building in the right structure.

\subsection{Outlook}

In this lecture, we introduced:

\begin{itemize}
    \item convolution as a structured linear operator,
    \item translation invariance,
    \item pooling and hierarchical representations,
    \item applications in vision.
\end{itemize}

\medskip

\noindent
Next, we study sequence models, which address data with temporal or sequential structure.

\medskip

\noindent
\textbf{Guiding principle.}  
Different data modalities require different architectural biases.